\section{Traces of relations}

\import{relation/properties.tex}

The trace of a relation restricts both coordinates to the displayed set.  It
is different from a domain restriction, which constrains only the first
coordinate.

\begin{definition}\label{relation_trace}
    $\relationTrace{R}{A}=R\inter(A\times A)$.
\end{definition}

\begin{proposition}\label{relation_trace_iff}
    $x\mathrel{\relationTrace{R}{A}}y$ iff
        $x\in A$ and $y\in A$ and $x\mathrel{R}y$.
\end{proposition}
\begin{proof}
    Follows by
        \auto{relation_trace,inter,times_tuple_elim,times_tuple_intro}.
\end{proof}

\begin{proposition}\label{relation_trace_subseteq}
    $\relationTrace{R}{A}\subseteq R$.
\end{proposition}
\begin{proof}
    Follows by \auto{relation_trace,inter_lower_left}.
\end{proof}

\begin{classification}\label{relation_trace_binary_on}
    $\relationTrace{R}{A}$ is a binary relation on $A$.
\end{classification}
\begin{proof}
    $\relationTrace{R}{A}\subseteq A\times A$ by
        \auto{relation_trace,inter_lower_right}.
    Follows by definition.
\end{proof}

\begin{proposition}\label{relation_trace_reflexive_on}
    Suppose $R$ is reflexive on $X$ and $A\subseteq X$.
    Then $\relationTrace{R}{A}$ is reflexive on $A$.
\end{proposition}
\begin{proof}
    Follows by
        \obv{reflexive_on,relation_trace_iff,elem_subseteq}.
\end{proof}

\begin{proposition}\label{relation_trace_transitive}
    Suppose $R$ is transitive.
    Then $\relationTrace{R}{A}$ is transitive.
\end{proposition}
\begin{proof}
    Follows by \auto{transitive,relation_trace_iff}.
\end{proof}

\begin{proposition}\label{relation_trace_antisymmetric}
    Suppose $R$ is antisymmetric.
    Then $\relationTrace{R}{A}$ is antisymmetric.
\end{proposition}
\begin{proof}
    Follows by \auto{antisymmetric,relation_trace_iff}.
\end{proof}

\begin{proposition}\label{relation_trace_symmetric}
    Suppose $R$ is symmetric.
    Then $\relationTrace{R}{A}$ is symmetric.
\end{proposition}
\begin{proof}
    Follows by \auto{symmetric,relation_trace_iff}.
\end{proof}

\begin{proposition}\label{relation_trace_irreflexive}
    Suppose $R$ is irreflexive.
    Then $\relationTrace{R}{A}$ is irreflexive.
\end{proposition}
\begin{proof}
    Follows by \auto{irreflexive,relation_trace_iff}.
\end{proof}

\begin{proposition}\label{relation_trace_asymmetric}
    Suppose $R$ is asymmetric.
    Then $\relationTrace{R}{A}$ is asymmetric.
\end{proposition}
\begin{proof}
    Follows by \auto{asymmetric,relation_trace_iff}.
\end{proof}

\begin{proposition}\label{relation_trace_connex_on}
    Suppose $R$ is connex on $X$ and $A\subseteq X$.
    Then $\relationTrace{R}{A}$ is connex on $A$.
\end{proposition}
\begin{proof}
    Follows by \auto{connex,relation_trace_iff,elem_subseteq}.
\end{proof}

\begin{proposition}\label{membership_relation_trace}
    Suppose $A\subseteq B$.
    Then $\relationTrace{\memrel{B}}{A}=\memrel{A}$.
\end{proposition}
\begin{proof}
    $\relationTrace{\memrel{B}}{A}$ is a relation by
        \auto{relation,relation_trace_subseteq,subseteq,memrel_is_relation}.
    $\memrel{A}$ is a relation by \obv{memrel_is_relation}.
    Follows by
        \auto{relext,relation_trace_iff,memrel,elem_subseteq}.
\end{proof}
